\section*{Chapitre \ref{ch:g6:density-floating} --- La masse volumique~: pourquoi les objets flottent}

\begin{solution}{exo:g6:density-floating:1}
La \omterm{def:g3:measuring-mass:mass}{masse} d'un \omterm{def:g6:mass-and-volume:volume}{centimètre cube} de la substance, en \omterm{def:g3:measuring-mass:units}{grammes}~; celle de
l'eau vaut \qty{1}{g} par \unit{cm^3}. La règle~: moins dense que
l'eau, ça \omterm{def:g1:floating-and-sinking:words}{flotte}~; plus dense, ça \omterm{def:g1:floating-and-sinking:words}{coule}.
\end{solution}

\begin{solution}{exo:g6:density-floating:2}
Le chêne \omterm{def:g1:floating-and-sinking:words}{flotte}~; la pierre \omterm{def:g1:floating-and-sinking:words}{coule}~; le liège \omterm{def:g1:floating-and-sinking:words}{flotte}~; le fer \omterm{def:g1:floating-and-sinking:words}{coule}~;
la glace \omterm{def:g1:floating-and-sinking:words}{flotte} --- d'un cheveu.
\end{solution}

\begin{solution}{exo:g6:density-floating:3}
$120 \div 80 = 1.5$~: \omterm{def:g6:density-floating:density}{masse volumique} $1.5$ --- plus dense que
l'eau, le bloc \omterm{def:g1:floating-and-sinking:words}{coule}.
\end{solution}

\begin{solution}{exo:g6:density-floating:4}
$158 \div 20 = 7.9$~: la signature du fer.
\end{solution}

\begin{solution}{exo:g6:density-floating:5}
La \omterm{def:g6:density-floating:density}{masse volumique} appartient à la substance, pas au morceau~: chaque
\omterm{def:g6:mass-and-volume:volume}{centimètre cube} de chêne porte la même part de \omterm{def:g3:measuring-mass:units}{grammes}, éclat ou
poutre. Cela fait de la \omterm{def:g6:density-floating:density}{masse volumique} une carte d'identité pour les
substances inconnues.
\end{solution}

\begin{solution}{exo:g6:density-floating:6}
Les \omterm{def:g3:solids-liquids-gases:liquid}{liquides} s'empilent par \omterm{def:g6:density-floating:density}{masse volumique}, le plus dense en bas~:
le miel $1.4$ sous l'eau $1.0$ sous l'huile $0.9$. Le grain de raisin
($1.05$) traverse l'huile et l'eau et se pose sur le plancher de
miel --- plus dense que les deux étages du dessus, moins dense que le
miel.
\end{solution}

\begin{solution}{exo:g6:density-floating:7}
La plupart des \omterm{def:g3:solids-liquids-gases:solid}{solides} sont plus denses que leur propre \omterm{def:g3:solids-liquids-gases:liquid}{liquide} et y
coulent. La glace \omterm{def:g1:floating-and-sinking:words}{flotte} parce que l'eau gonfle en gelant --- la même
matière dans un dixième de place en plus donne une \omterm{def:g6:density-floating:density}{masse volumique}
$0.9 < 1$.
\end{solution}

\begin{solution}{exo:g6:density-floating:8}
Le bois signe environ $0.6$~: chaque \omterm{def:g6:mass-and-volume:volume}{centimètre cube} du tronc fait
une offre inférieure au \qty{1}{g} de l'eau, et le tronc \omterm{def:g1:floating-and-sinking:words}{flotte} donc,
aussi énorme soit-il. Le bronze signe près de $9$~: chaque \omterm{def:g6:mass-and-volume:volume}{centimètre
cube} de la pièce surenchérit sur l'eau, et elle \omterm{def:g1:floating-and-sinking:words}{coule} donc, aussi
petite soit-elle. La taille n'a jamais compté~; la \omterm{def:g3:measuring-mass:mass}{masse} par \omterm{def:g6:measurement-in-science:measuring}{unité} de
place a toujours compté.
\end{solution}

\begin{solution}{exo:g6:density-floating:9}
$\qty{1}{L} = \qty{1000}{cm^3}$, donc $800 \div 1000 = 0.8$ \omterm{def:g3:measuring-mass:units}{gramme}
par \omterm{def:g6:mass-and-volume:volume}{centimètre cube}. La glace, à $0.9$, est \emph{plus dense} que ce
\omterm{def:g3:solids-liquids-gases:liquid}{liquide}~: le glaçon y \omterm{def:g1:floating-and-sinking:words}{coule}.
\end{solution}

\begin{solution}{exo:g6:density-floating:10}
Le colis flottant est la coque tout entière~: la coquille d'acier
\emph{plus} le grand \omterm{def:g6:mass-and-volume:volume}{volume} d'\omterm{def:g2:air-around-us:air}{air} qu'elle enferme. Moyennée sur toute
cette place, la \omterm{def:g6:density-floating:density}{masse volumique} du colis passe au-dessous de $1$, et
la règle dit~: ça \omterm{def:g1:floating-and-sinking:words}{flotte}. Une brèche laisse l'eau remplacer l'\omterm{def:g2:air-around-us:air}{air}~: la
\omterm{def:g6:density-floating:density}{masse volumique} du colis grimpe vers celle de l'acier, dépasse $1$
--- et la règle, imperturbable, dit~: ça \omterm{def:g1:floating-and-sinking:words}{coule}.
\end{solution}

\begin{solution}{exo:g6:density-floating:11}
L'arbitre a changé~: face au $1.03$ de l'eau de mer, le corps d'un
nageur, proche de $1$, se trouve confortablement du côté des
flotteurs, et la mer le porte donc plus haut. En entrant dans l'eau
douce ($1.0$), le navire perd cette marge et s'enfonce davantage~: sa
ligne de flottaison remonte sur la coque.
\end{solution}

\begin{solution}{exo:g6:density-floating:12}
À moitié immergée, la sphère s'enfonce exactement comme doit le faire
un colis de \omterm{def:g6:density-floating:density}{masse volumique} égale à la moitié de celle de l'eau~: la
coquille plus l'\omterm{def:g2:air-around-us:air}{air} donnent en moyenne environ $0.5$ \omterm{def:g3:measuring-mass:units}{gramme} par
\omterm{def:g6:mass-and-volume:volume}{centimètre cube} --- c'est l'\omterm{def:g2:air-around-us:air}{air} caché qui allège le plus. À mesure
que l'eau s'infiltre, elle remplace l'\omterm{def:g2:air-around-us:air}{air}, la \omterm{def:g6:density-floating:density}{masse volumique} moyenne
du colis grimpe, et la sphère s'enfonce de plus en plus~; à l'instant
où cette moyenne dépasse le $1.0$ de l'eau, la règle de la flottaison
change de camp et la sphère descend.
\end{solution}

\begin{solution}{pb:g6:density-floating:1}
\textbf{1.} Rien~: échanger de l'or contre une \emph{\omterm{def:g3:measuring-mass:mass}{masse}} égale
d'argent laisse la \omterm{def:g3:levers-and-scales:scale}{balance} satisfaite --- la \omterm{def:g3:measuring-mass:mass}{masse} seule ne peut pas
voir la substitution.
\textbf{2.} $1930 \div 19.3 = \qty{100}{cm^3}$.
\textbf{3.} \qty{130}{cm^3} --- le débordement est le \omterm{def:g6:mass-and-volume:volume}{volume} de la
couronne, par déplacement.
\textbf{4.} La couronne occupe \qty{30}{cm^3} de plus que ne le
devrait de l'or pur de cette \omterm{def:g3:measuring-mass:mass}{masse}~: quelque chose de plus volumineux
que l'or se cache à l'intérieur.
\textbf{5.} $1930 \div 130 \approx 14.8$.
\textbf{6.} L'or $19.3$ --- la couronne $14.8$ --- l'argent $10.5$~:
la couronne tombe entre les deux signatures.
\textbf{7.} «~Sire, chaque \omterm{def:g6:mass-and-volume:volume}{centimètre cube} d'or pur porte
\qty{19.3}{g}~; votre couronne n'en porte qu'environ $15$ --- de l'or
pur, elle ne l'est pas.~»
\textbf{8.} L'intervalle donne des \omterm{def:g6:density-floating:density}{masses volumiques} de
$1900 \div 130 \approx 14.6$ à $1960 \div 130 \approx 15.1$ --- très
loin de $19.3$~: aucune erreur de \omterm{def:g3:levers-and-scales:scale}{balance} plausible ne sauve
l'orfèvre.
\textbf{9.} L'argent est moins dense~: chaque \omterm{def:g3:measuring-mass:units}{gramme} d'or volé a été
rendu sous forme d'un \omterm{def:g3:measuring-mass:units}{gramme} d'argent qui exige \emph{plus de
place} --- une substitution \omterm{def:g3:measuring-mass:units}{gramme} pour \omterm{def:g3:measuring-mass:units}{gramme} conserve la \omterm{def:g3:measuring-mass:mass}{masse} et
gonfle le \omterm{def:g6:mass-and-volume:volume}{volume}.
\textbf{10.} Part d'or~: $965 \div 19.3 = \qty{50.0}{cm^3}$~; part
d'argent~: $965 \div 10.5 \approx \qty{91.9}{cm^3}$.
\textbf{11.} Total $\approx \qty{142}{cm^3}$ --- plus que les $130$
mesurés~: l'hypothèse moitié-moitié met trop d'argent~; le vrai vol
était plus modeste (le \omterm{def:g6:mass-and-volume:volume}{volume} mesuré se situe entre $100$ et $142$).
\textbf{12.} «~Prouvé~: la couronne n'est pas en or pur. Méthode~: sa
\omterm{def:g3:measuring-mass:mass}{masse} partagée sur son \omterm{def:g6:mass-and-volume:volume}{volume} mesuré par l'eau donne une \omterm{def:g6:density-floating:density}{masse
volumique} très inférieure à l'immuable signature de l'or. Et la
couronne n'a pas eu besoin de la moindre égratignure~: l'eau a lu son
secret de l'extérieur --- voilà le triomphe.~»
\end{solution}
